\documentclass{amsart}
% For dealing with references we use the comment environment
\usepackage{verbatim}
\newenvironment{reference}{\comment}{\endcomment}
%\newenvironment{reference}{}{}
\newenvironment{slogan}{\comment}{\endcomment}
\newenvironment{history}{\comment}{\endcomment}

% For commutative diagrams we use Xy-pic
\usepackage[all]{xy}

% We use 2cell for 2-commutative diagrams.
\xyoption{2cell}
\UseAllTwocells

% We use multicol for the list of chapters between chapters
\usepackage{multicol}

% This is generally recommended for better output
\usepackage{lmodern}
\usepackage[T1]{fontenc}

% For cross-file-references

% Package for hypertext links:
\usepackage{hyperref}

% For any local file, say "hello.tex" you want to link to please

% Theorem environments.
%
\theoremstyle{plain}
\newtheorem{theorem}[subsection]{Theorem}
\newtheorem{proposition}[subsection]{Proposition}
\newtheorem{lemma}[subsection]{Lemma}

\theoremstyle{definition}
\newtheorem{definition}[subsection]{Definition}
\newtheorem{example}[subsection]{Example}
\newtheorem{exercise}[subsection]{Exercise}
\newtheorem{situation}[subsection]{Situation}

\theoremstyle{remark}
\newtheorem{remark}[subsection]{Remark}
\newtheorem{remarks}[subsection]{Remarks}

\numberwithin{equation}{subsection}

% Macros
%
\def\lim{\mathop{\mathrm{lim}}\nolimits}
\def\colim{\mathop{\mathrm{colim}}\nolimits}
\def\Spec{\mathop{\mathrm{Spec}}}
\def\Hom{\mathop{\mathrm{Hom}}\nolimits}
\def\Ext{\mathop{\mathrm{Ext}}\nolimits}
\def\SheafHom{\mathop{\mathcal{H}\!\mathit{om}}\nolimits}
\def\SheafExt{\mathop{\mathcal{E}\!\mathit{xt}}\nolimits}
\def\Sch{\mathit{Sch}}
\def\Mor{\mathop{\mathrm{Mor}}\nolimits}
\def\Ob{\mathop{\mathrm{Ob}}\nolimits}
\def\Sh{\mathop{\mathit{Sh}}\nolimits}
\def\NL{\mathop{N\!L}\nolimits}
\def\CH{\mathop{\mathrm{CH}}\nolimits}
\def\proetale{{pro\text{-}\acute{e}tale}}
\def\etale{{\acute{e}tale}}
\def\QCoh{\mathit{QCoh}}
\def\Ker{\mathop{\mathrm{Ker}}}
\def\Im{\mathop{\mathrm{Im}}}
\def\Coker{\mathop{\mathrm{Coker}}}
\def\Coim{\mathop{\mathrm{Coim}}}

% Boxtimes
%
\DeclareMathSymbol{\boxtimes}{\mathbin}{AMSa}{"02}

%
% Macros for moduli stacks/spaces
%
\def\QCohstack{\mathcal{QC}\!\mathit{oh}}
\def\Cohstack{\mathcal{C}\!\mathit{oh}}
\def\Spacesstack{\mathcal{S}\!\mathit{paces}}
\def\Quotfunctor{\mathrm{Quot}}
\def\Hilbfunctor{\mathrm{Hilb}}
\def\Curvesstack{\mathcal{C}\!\mathit{urves}}
\def\Polarizedstack{\mathcal{P}\!\mathit{olarized}}
\def\Complexesstack{\mathcal{C}\!\mathit{omplexes}}
% \Pic is the operator that assigns to X its picard group, usage \Pic(X)
% \Picardstack_{X/B} denotes the Picard stack of X over B
% \Picardfunctor_{X/B} denotes the Picard functor of X over B
\def\Pic{\mathop{\mathrm{Pic}}\nolimits}
\def\Picardstack{\mathcal{P}\!\mathit{ic}}
\def\Picardfunctor{\mathrm{Pic}}
\def\Deformationcategory{\mathcal{D}\!\mathit{ef}}

\setlength{\emergencystretch}{3em}
\begin{document}
\title[Stacks Project, Tag 07M8]{727 --- Smooth and syntomic algebras lift along arbitrary ring surjections}
\maketitle
\noindent Copyright (C) 2005--2025 Johan de Jong. Stacks Project authors. Source: \href{https://stacks.math.columbia.edu/tag/07M8}{Tag 07M8}.

\noindent Editorial difficulty note (not author text). Difficulty H2: The proof improves a conormal presentation, lifts a complementary projective module through an etale extension, adds a vector bundle, and kills lifted polynomial coordinates. This constructive lifting theorem requires substantial smoothing and cotangent-complex machinery..

\noindent Editorial provenance note (not author text). History: excerpt from revision \texttt{a0444\allowbreak 6e57e\allowbreak c1fbc\allowbreak 252a8\allowbreak 71afc\allowbreak ec775\allowbreak 2fb28\allowbreak 07b14}, smoothing.tex, lines 447--557. Reference presentation was adapted; original-author context and core support blocks were retained or added. The mathematical wording is unchanged. Licensed under GNU FDL 1.2 or later; no invariant sections or cover texts. See ../licenses/COPYING. Context blocks reproduce author definitions and supporting results; they are not additional counted problems.

\begin{definition}
\label{more-algebra-definition-local-complete-intersection}
A ring map $A \to B$ is called a {\it local complete intersection}
if it is of finite type and for some (equivalently any) presentation
$B = A[x_1, \ldots, x_n]/I$ the ideal $I$ is Koszul-regular.
\end{definition}

\begin{lemma}
\label{lemma-improve-presentation}
Let $R$ be a ring and let $A$ be a finitely presented $R$-algebra.
There exists finite type $R$-algebra map $A \to C$ which has a
retraction with the following two properties
\begin{enumerate}
\item for each $a \in A$ such that $R \to A_a$ is a local complete
intersection (More on Algebra, Definition
\ref{more-algebra-definition-local-complete-intersection})
the ring $C_a$ is smooth over $A_a$ and has a presentation
$C_a = R[y_1, \ldots, y_m]/J$ such that $J/J^2$ is free over $C_a$, and
\item for each $a \in A$ such that $A_a$ is smooth over $R$ the
module $\Omega_{C_a/R}$ is free over $C_a$.
\end{enumerate}
\end{lemma}

\begin{proof}
Choose a presentation $A = R[x_1, \ldots, x_n]/I$ and write
$I = (f_1, \ldots, f_m)$. Define the $A$-module $K$ by the short exact sequence
$$
0 \to K \to A^{\oplus m} \to I/I^2 \to 0
$$
where the $j$th basis vector $e_j$ in the middle is mapped to the class of
$f_j$ on the right. Set
$$
C = \text{Sym}^*_A(I/I^2).
$$
The retraction is just the projection onto the degree $0$ part of $C$.
We have a surjection $R[x_1, \ldots, x_n, y_1, \ldots, y_m] \to C$
which maps $y_j$ to the class of $f_j$ in $I/I^2$. The kernel $J$ of this
map is generated by the elements $f_1, \ldots, f_m$ and by elements
$\sum h_j y_j$ with $h_j \in R[x_1, \ldots, x_n]$ such that
$\sum h_j e_j$ defines an element of $K$. By
Algebra, Lemma \href{https://stacks.math.columbia.edu/tag/00S2}{lemma exact sequence NL (Tag 00S2)}
applied to $R \to A \to C$ and the presentations above and
More on Algebra, Lemma
\href{https://stacks.math.columbia.edu/tag/07EV}{lemma cotangent complex symmetric algebra (Tag 07EV)}
there is a short exact sequence
\begin{equation}
\label{equation-sequence}
I/I^2 \otimes_A C \to J/J^2 \to K \otimes_A C \to 0
\end{equation}
of $C$-modules. Let $h \in R[x_1, \ldots, x_n]$ be an element
with image $a \in A$. We will use as presentations for the localized rings
$$
A_a = R[x_0, x_1, \ldots, x_n]/I'
\quad\text{and}\quad
C_a = R[x_0, x_1, \ldots, x_n, y_1, \ldots, y_m]/J'
$$
where $I' = (hx_0 - 1, I)$ and $J' = (hx_0 - 1, J)$. Hence
$I'/(I')^2 = A_a \oplus (I/I^2)_a$ as $A_a$-modules and
$J'/(J')^2 = C_a \oplus (J/J^2)_a$ as $C_a$-modules.
Thus we obtain
\begin{equation}
\label{equation-sequence-localized}
C_a \oplus I/I^2 \otimes_A C_a \to
C_a \oplus (J/J^2)_a \to
K \otimes_A C_a \to 0
\end{equation}
as the sequence of
Algebra, Lemma \href{https://stacks.math.columbia.edu/tag/00S2}{lemma exact sequence NL (Tag 00S2)}
corresponding to $R \to A_a \to C_a$ and the presentations above.

\medskip\noindent
Next, assume that $a \in A$ is such that $A_a$ is a local complete
intersection over $R$. Then $(I/I^2)_a$ is finite projective over $A_a$, see
More on Algebra, Lemma
\href{https://stacks.math.columbia.edu/tag/08RK}{lemma quasi regular ideal finite projective (Tag 08RK)}.
Hence we see $K_a \oplus (I/I^2)_a \cong A_a^{\oplus m}$ is free.
In particular $K_a$ is finite projective too.
By More on Algebra, Lemma \href{https://stacks.math.columbia.edu/tag/07D4}{lemma transitive lci at end (Tag 07D4)}
the sequence (\ref{equation-sequence-localized}) is exact on the left.
Hence
$$
J'/(J')^2 \cong
C_a \oplus I/I^2 \otimes_A C_a \oplus K \otimes_A C_a \cong
C_a^{\oplus m + 1}
$$
This proves (1). Finally, suppose that in addition $A_a$ is smooth over
$R$. Then the same presentation shows that $\Omega_{C_a/R}$
is the cokernel of the map
$$
J'/(J')^2 \longrightarrow
\bigoplus\nolimits_i C_a\text{d}x_i \oplus \bigoplus\nolimits_j C_a\text{d}y_j
$$
The summand $C_a$ of $J'/(J')^2$ in the decomposition above
corresponds to $hx_0 - 1$ and hence maps
isomorphically to the summand $C_a\text{d}x_0$. The summand
$I/I^2 \otimes_A C_a$ of $J'/(J')^2$ maps injectively to
$\bigoplus_{i = 1, \ldots, n} C_a\text{d}x_i$
with quotient $\Omega_{A_a/R} \otimes_{A_a} C_a$. The summand
$K \otimes_A C_a$ maps injectively to
$\bigoplus_{j \geq 1} C_a\text{d}y_j$ with quotient isomorphic to
$I/I^2 \otimes_A C_a$. Thus the cokernel of the last displayed
map is the module
$I/I^2 \otimes_A C_a \oplus \Omega_{A_a/R} \otimes_{A_a} C_a$.
Since $(I/I^2)_a \oplus \Omega_{A_a/R}$ is
free (from the definition of smooth ring maps) we see that (2) holds.
\end{proof}


\begin{proposition}
\label{proposition-lift-smooth}
\begin{slogan}
Smooth and syntomic algebras lift along surjections
\end{slogan}
Let $R \to R_0$ be a surjective ring map with kernel $I$.
\begin{enumerate}
\item If $R_0 \to A_0$ is a syntomic ring map, then there exists a syntomic
ring map $R \to A$ such that $A/IA \cong A_0$.
\item If $R_0 \to A_0$ is a smooth ring map, then there exists a smooth
ring map $R \to A$ such that $A/IA \cong A_0$.
\end{enumerate}
\end{proposition}

\begin{proof}
Assume $R_0 \to A_0$ syntomic, in particular a local complete intersection
(More on Algebra, Lemma \href{https://stacks.math.columbia.edu/tag/07D3}{lemma syntomic lci (Tag 07D3)}).
Choose a presentation $A_0 = R_0[x_1, \ldots, x_n]/J_0$. Set
$C_0 = \text{Sym}^*_{A_0}(J_0/J_0^2)$. Note that $J_0/J_0^2$ is a finite
projective $A_0$-module (Algebra, Lemma
\href{https://stacks.math.columbia.edu/tag/07BT}{lemma syntomic presentation ideal mod squares (Tag 07BT)}).
By Lemma \ref{lemma-improve-presentation} the ring map
$A_0 \to C_0$ is smooth and we can find a presentation
$C_0 = R_0[y_1, \ldots, y_m]/K_0$ with $K_0/K_0^2$ free over $C_0$.
By Algebra, Lemma \href{https://stacks.math.columbia.edu/tag/07CF}{lemma huber (Tag 07CF)} we can assume
$C_0 = R_0[y_1, \ldots, y_m]/(\overline{f}_1, \ldots, \overline{f}_c)$
where $\overline{f}_1, \ldots, \overline{f}_c$ maps to a basis of
$K_0/K_0^2$ over $C_0$. Choose
$f_1, \ldots, f_c \in R[y_1, \ldots, y_c]$ lifting
$\overline{f}_1, \ldots, \overline{f}_c$ and set
$$
C = R[y_1, \ldots, y_m]/(f_1, \ldots, f_c)
$$
By construction $C_0 = C/IC$. By Algebra, Lemma
\href{https://stacks.math.columbia.edu/tag/00ST}{lemma localize relative complete intersection (Tag 00ST)}
we can after replacing $C$ by $C_g$ assume that $C$ is a relative
global complete intersection over $R$.
We conclude that there exists a finite projective $A_0$-module
$P_0$ such that $C_0 = \text{Sym}^*_{A_0}(P_0)$
is isomorphic to $C/IC$ for some syntomic $R$-algebra $C$.

\medskip\noindent
Choose an integer $n$ and a direct sum decomposition
$A_0^{\oplus n} = P_0 \oplus Q_0$.
By More on Algebra, Lemma \href{https://stacks.math.columbia.edu/tag/07M5}{lemma lift projective module (Tag 07M5)}
we can find an \'etale ring map $C \to C'$ which induces
an isomorphism $C/IC \to C'/IC'$ and a finite projective
$C'$-module $Q$ such that $Q/IQ$ is isomorphic to
$Q_0 \otimes_{A_0} C/IC$.
Then $D = \text{Sym}_{C'}^*(Q)$ is a smooth $C'$-algebra (see
More on Algebra, Lemma \href{https://stacks.math.columbia.edu/tag/07M6}{lemma symmetric algebra smooth (Tag 07M6)}).
Picture
$$
\xymatrix{
R \ar[d] \ar[rr] & &
C \ar[r] \ar[d] &
C' \ar[r] \ar[d] &
D \ar[d] \\
R/I \ar[r] &
A_0 \ar[r] &
C/IC \ar[r]^{\cong} &
C'/IC' \ar[r] &
D/ID
}
$$
Observe that our choice of $Q$ gives
\begin{align*}
D/ID & =
\text{Sym}_{C/IC}^*(Q_0 \otimes_{A_0} C/IC) \\
& =
\text{Sym}_{A_0}^*(Q_0) \otimes_{A_0} C/IC \\
& =
\text{Sym}_{A_0}^*(Q_0) \otimes_{A_0}
\text{Sym}_{A_0}^*(P_0) \\
& =
\text{Sym}_{A_0}^*(Q_0 \oplus P_0) \\
& =
\text{Sym}_{A_0}^*(A_0^{\oplus n}) \\
& =
A_0[x_1, \ldots, x_n]
\end{align*}
Choose $f_1, \ldots, f_n \in D$ which map to $x_1, \ldots, x_n$
in $D/ID = A_0[x_1, \ldots, x_n]$. Set $A = D/(f_1, \ldots, f_n)$.
Note that $A_0 = A/IA$. We claim that $R \to A$ is syntomic
in a neighbourhood of $V(IA)$. If the claim is true, then we can
find a $f \in A$ mapping to $1 \in A_0$ such that $A_f$ is syntomic
over $R$ and the proof of (1) is finished.

\medskip\noindent
Proof of the claim. Observe that $R \to D$ is syntomic as a composition
of the syntomic ring map $R \to C$, the \'etale ring map $C \to C'$ and the
smooth ring map $C' \to D$ (Algebra, Lemmas
\href{https://stacks.math.columbia.edu/tag/00SZ}{lemma composition syntomic (Tag 00SZ)} and
\href{https://stacks.math.columbia.edu/tag/00TA}{lemma smooth syntomic (Tag 00TA)}).
The question is local on $\Spec(D)$, hence we
may assume that $D$ is a relative global complete intersection
(Algebra, Lemma \href{https://stacks.math.columbia.edu/tag/00SY}{lemma syntomic (Tag 00SY)}).
Say $D = R[y_1, \ldots, y_m]/(g_1, \ldots, g_s)$.
Let $f'_1, \ldots, f'_n \in R[y_1, \ldots, y_m]$ be lifts of
$f_1, \ldots, f_n$. Then we can apply
Algebra, Lemma \href{https://stacks.math.columbia.edu/tag/00ST}{lemma localize relative complete intersection (Tag 00ST)}
to get the claim.

\medskip\noindent
Proof of (2). Since a smooth ring map is syntomic, we can find
a syntomic ring map $R \to A$ such that $A_0 = A/IA$.
By assumption the fibres of $R \to A$ are smooth over primes in $V(I)$
hence $R \to A$ is smooth in an open neighbourhood of $V(IA)$
(Algebra, Lemma \href{https://stacks.math.columbia.edu/tag/00TF}{lemma flat fibre smooth (Tag 00TF)}).
Thus we can replace $A$ by a localization to obtain the result we want.
\end{proof}
\end{document}
